\documentclass[12pt,a4paper]{article}

\usepackage{fontawesome5}
\providecommand{\faPencilAlt}{\faIcon{pencil-alt}}

\usepackage{styles/couleurs}
\usepackage{styles/mise_en_page}
\usepackage{styles/boites}
\usepackage{styles/maths}
\usepackage{styles/figures}
\usepackage{styles/evaluations}

\usepackage{enumitem}
\usepackage{array}
\usepackage{tabularx}
\usepackage{amsmath}

\setlist[itemize]{leftmargin=1.4em,itemsep=0.15em,topsep=0.2em}
\setlist[enumerate]{leftmargin=*,itemsep=0.25em,topsep=0.25em}

\begin{document}

\enteteinterro[Calculatrice interdite sauf mention contraire]{4e}{Quotients et enchaînements}{15 min}{12}

\exo{3}[Calculer]
Calculer les quotients. \pts{3}

\[
(-42)\div(+6)
\qquad
(-35)\div(-5)
\qquad
7{,}2\div(-4)
\]

\lignes{4}

\exo{2}[Raisonner]
Déterminer le signe du produit, puis calculer. \pts{2}

\[
(-2)\times3\times(-5)\times(-1)
\]

\lignes{3}

\exo{4}[Calculer]
Calculer en détaillant les étapes. \pts{4}

\[
A=(-24)\div(-6)+5\times(-2)
\]

\[
B=7-(-4)\times(-3)
\]

\lignes{7}

\exo{3}[Raisonner]
Un élève écrit :

\[
18-(-3)\times4=(18-(-3))\times4=21\times4=84.
\]

Expliquer l'erreur et corriger le calcul. \pts{3}

\lignes{5}

\end{document}
